\section{Computational Experiments}
\label{sec:experiments}

This section evaluates the proposed SND-RTI model through a series of
computational experiments on randomly generated instances that approximate
real-world manufacturing supply chains. The experiments pursue two
objectives: (i)~assessing the computational scalability of the linearised
MILP as the network grows; and (ii)~quantifying the cost savings of the
optimised collaborative return strategy relative to the conventional
one-to-one baseline.

% ══════════════════════════════════════════════════════════════════
%  INSTANCE GENERATION
% ══════════════════════════════════════════════════════════════════

\subsection{Instance Generation}
\label{sec:instance-gen}

Instances are generated to mirror the structure of multi-tier
manufacturing supply chains, where plants are geographically clustered
in industrial zones and connected by transport routes of varying
distance and mode availability.

\paragraph{Network topology.}
Each instance is built in three steps.

\begin{enumerate}
  \item \textbf{Industrial zones.}
        $n_Z$ industrial zones are placed uniformly at random in a
        $[0, L] \times [0, L]$~km\textsuperscript{2} area, where
        $L$ is chosen according to the instance class
        (Table~\ref{tab:instance-classes}).  Each zone~$z$ contains
        $n^{L}_z$ large plants and $n^{S}_z$ small plants, drawn
        from the ranges specified in Table~\ref{tab:instance-classes}.
        Plants within a zone are scattered around the zone centre
        with a Gaussian offset ($\sigma_{\text{intra}}$~km).

  \item \textbf{Route generation.}
        A full-RTI route $(i,j)$ is generated between every
        ordered pair of plants that satisfies a demand condition
        (see below). Large plants have a higher probability
        $p^{\text{inter}}$ of sourcing from plants in \emph{other}
        zones, while small plants have a higher probability
        $p^{\text{intra}}$ of sourcing within the \emph{same}
        zone. Specifically, for each plant~$j$, a set of suppliers
        is sampled: with probability $p^{\text{inter}}$ (resp.\
        $p^{\text{intra}}$) a candidate supplier is drawn from
        another zone (resp.\ the same zone). The expected number of
        supply routes per plant is controlled by parameter
        $\bar{k}$. Empty-RTI return routes are allowed between any
        pair of nodes in the network, i.e.\ $be_{ijmr}=1$ for all
        $(i,j) \in V \times V$, $m \in \mathcal{M}$, $r \in R$,
        since the model optimises which of these potential return
        arcs are actually used.

  \item \textbf{Hub designation.}
        Each industrial zone contains $n_H$ candidate hub nodes
        ($H \subseteq V$), reflecting the common practice of
        co-locating consolidation centres within industrial areas.
        The value of $n_H$ scales with the instance class:
        $n_H = 1$ for \textsc{Small}, $n_H = 2$ for
        \textsc{Medium}, and $n_H = 3$ for \textsc{Large}
        instances (see Table~\ref{tab:instance-classes}). Within
        each zone, hub candidates are selected among the nodes
        with the highest total inflow plus outflow of full RTIs.
        The total number of candidate hubs is therefore
        $|H| = n_H \cdot n_Z$; the model decides which of these
        candidates to activate.
\end{enumerate}

\paragraph{Demand and product assignment.}
Each full-RTI route $(i,j)$ carries one or more products from the
product set~$P$. The number of products per route depends on the
instance class and the type of the receiving plant: routes serving
large plants carry up to~2 products in \textsc{Small} instances,
up to~3 in \textsc{Medium}, and up to~4 in \textsc{Large}
instances; routes serving small plants carry a single product. For
each product assigned to a route, the daily average demand in full
RTI units, $\mu_{ijmpr}$, is drawn from a log-normal distribution
$\text{LogNormal}(\mu_d, \sigma_d)$, rounded to the nearest integer,
to reflect the right-skewed demand patterns typical of manufacturing
supply chains. Large plants receive demands from the upper quartile
of the distribution to model their higher throughput. The inlay
fraction $ric_{ijpr}$ is drawn uniformly from $[0, r^{\max}]$.

\paragraph{RTI types and transport modes.}
Each product~$p$ is compatible with a subset of RTI types
$R_p \subseteq R$. The majority of products are assigned to a
single RTI type ($|R_p|=1$), chosen uniformly at random from~$R$.
A smaller fraction of products admits multiple compatible types:
up to~2 in \textsc{Small} instances, 2--3 in \textsc{Medium}, and
up to~4 in \textsc{Large} instances. The parameter $bf_{ijmrp}$ is
set to~1 for every compatible product--RTI-type pair on routes where
product~$p$ is demanded; the model selects the type to use.
Two transport modes are considered throughout:
$\mathcal{M} = \{\text{FTL}, \text{LTL}\}$. Mode availability
$bm_{ijm}$ depends on distance: routes shorter than
$d^{\text{FTL}}_{\min}$~km are LTL-only; routes longer than
$d^{\text{LTL}}_{\max}$~km are FTL-only; all others admit both modes.

\paragraph{Cost and capacity parameters.}
Variable transport costs $c_{ijm}$ are set proportional to the
Euclidean distance between nodes, with mode-specific rates
$c^{\text{FTL}}$ and $c^{\text{LTL}}$ per km per m\textsuperscript{3}.
Fixed dispatch costs $f_{ijm}$ are mode-dependent constants.
Vehicle capacities $\delta_{ijm}$ and minimum load factors
$\omega_{ijm}$ follow industry-standard values for full-truckload
and groupage shipments. Transit times $l_{ijm}$ are obtained by
dividing the Euclidean distance by an average speed of
$v_{\text{avg}} = 500$~km/day, common to both modes.
Mode availability thresholds are set to
$d^{\text{FTL}}_{\min} = 50$~km and $d^{\text{LTL}}_{\max} = 800$~km.

\paragraph{Safety stocks and existing pool.}
Safety stocks at both endpoints of each route are set to three days
of the daily average demand on that route: $ssf^{i}_{ijpr} =
ssf^{j}_{ijpr} = 3\,\mu_{ijmpr}$ for full RTIs, and
$sse^{i}_{ijpr} = sse^{j}_{ijpr} = 3\,NE_{ijmr}$ for empty RTIs.
The existing RTI pool is set to zero for all types
($s^{\text{RTI}}_r = 0$), so that the model must purchase the
entire pool; this isolates the effect of the network design on
total RTI investment. Non-zero initial pools are reserved for the
case study, where historical inventory data are available.

\paragraph{Linearisation.}
The continuous queue time $q_{ijm}$ is discretised over the set of
admissible integer values
$\mathcal{Q}_m = \{q^{\min}_m, q^{\min}_m{+}1, \ldots,
q^{\max}_m\}$, yielding $|\mathcal{Q}_{\text{FTL}}| = 8$ and
$|\mathcal{Q}_{\text{LTL}}| = 5$ candidate values. Each
candidate~$k \in \mathcal{Q}_m$ is associated with a binary
variable indicating whether it is selected; all bilinear products
of the form $q_{ijm} \cdot (\cdot)$ are replaced by their
linearised counterparts as described in the model formulation.

RTI volumes $v^f_r$ and $v^e_r$, purchase costs $p^{\text{RTI}}_r$,
loss factors $\alpha^f$ and $\alpha^e$, and hub costs $fhc_i$ and
$vhc_i$ are sampled from the ranges reported in
Table~\ref{tab:param-ranges}.

% ══════════════════════════════════════════════════════════════════
%  INSTANCE CLASSES
% ══════════════════════════════════════════════════════════════════

\subsection{Instance Classes}
\label{sec:instance-classes}

Three classes of instances---\textsc{Small}, \textsc{Medium}, and
\textsc{Large}---are defined by scaling the number of industrial
zones, plants per zone, RTI types, and products.
Table~\ref{tab:instance-classes} summarises the dimensions.
For each class, $N_{\text{inst}}$ instances are generated with
different random seeds to allow statistical comparison.

\begin{table}[ht]
\centering
\caption{Instance classes and their dimensions.}
\label{tab:instance-classes}
\small
\begin{tabular}{lrrrrrrrrrrr}
\toprule
\textbf{Class}
  & $n_Z$ & $n^L_z$ & $n^S_z$ & $n_H$ & $|V|$ & $|E|$ (approx.)
  & $|R|$ & $\max|R_p|$ & $|P|$ & $|\mathcal{M}|$ & $N_{\text{inst}}$ \\
\midrule
\textsc{Small}
  & 2--3  & 1--2  & 2--4  & 1 & 10--20  & 15--40
  & 2     & 2  & 3--5  & 2 & 10 \\
\textsc{Medium}
  & 4--6  & 2--3  & 4--8  & 2 & 30--60  & 60--150
  & 3--4  & 3  & 6--10 & 2 & 10 \\
\textsc{Large}
  & 8--12 & 3--5  & 6--12 & 3 & 80--180 & 200--600
  & 4--6  & 4  & 10--20 & 2 & 10 \\
\bottomrule
\end{tabular}
\end{table}

\begin{table}[ht]
\centering
\caption{Parameter ranges for instance generation.}
\label{tab:param-ranges}
\small
\begin{tabular}{lll}
\toprule
\textbf{Parameter} & \textbf{Range / Value} & \textbf{Unit} \\
\midrule
\multicolumn{3}{l}{\textit{Network generation}} \\
$L$ (area side)
  & 500 / 1\,000 / 1\,500  & km \\
$\sigma_{\text{intra}}$
  & 5--15  & km \\
$\bar{k}$ (avg.\ suppliers per plant)
  & 2--4  & --- \\
$p^{\text{inter}}$ (large plants)
  & 0.7  & --- \\
$p^{\text{intra}}$ (small plants)
  & 0.7  & --- \\
$n_H$ (candidate hubs per zone)
  & 1 / 2 / 3  & --- \\
\midrule
\multicolumn{3}{l}{\textit{Demand and products}} \\
$\mu_d, \sigma_d$ (demand log-normal)
  & (2.0,\,0.8)  & RTI/day \\
$r^{\max}$ (max inlay fraction)
  & 0.3  & --- \\
$|R_p|$ (compatible RTI types per product)
  & 1 (majority); up to 2\,/\,3\,/\,4  & --- \\
\midrule
\multicolumn{3}{l}{\textit{Transport}} \\
$c^{\text{FTL}},\; c^{\text{LTL}}$
  & 0.05,\; 0.12  & \euro/km/m\textsuperscript{3} \\
$f^{\text{FTL}},\; f^{\text{LTL}}$
  & 150,\; 40  & \euro/shipment \\
$\delta^{\text{FTL}},\; \delta^{\text{LTL}}$
  & 80,\; 30  & m\textsuperscript{3} \\
$\omega^{\text{FTL}},\; \omega^{\text{LTL}}$
  & 40,\; 5  & m\textsuperscript{3} \\
$v_{\text{avg}}$ (average speed)
  & 500  & km/day \\
$d^{\text{FTL}}_{\min}$ (FTL minimum distance)
  & 50  & km \\
$d^{\text{LTL}}_{\max}$ (LTL maximum distance)
  & 800  & km \\
$[q^{\min},q^{\max}]_{\text{FTL}}$
  & [3,\,10]  & days \\
$[q^{\min},q^{\max}]_{\text{LTL}}$
  & [1,\,5]  & days \\
$|\mathcal{Q}_{\text{FTL}}|,\; |\mathcal{Q}_{\text{LTL}}|$
  & 8,\; 5  & --- \\
\midrule
\multicolumn{3}{l}{\textit{RTI and inventory}} \\
$v^f_r$
  & 0.3--0.8  & m\textsuperscript{3} \\
$v^e_r / v^f_r$ (folding ratio)
  & 0.2--0.5  & --- \\
$p^{\text{RTI}}_r$
  & 30--120  & \euro \\
$s^{\text{RTI}}_r$ (existing pool)
  & 0  & --- \\
$ss_{\text{days}}$ (safety stock)
  & 3  & days \\
$\alpha^f,\; \alpha^e$
  & 0.02--0.05  & --- \\
\midrule
\multicolumn{3}{l}{\textit{Hubs}} \\
$fhc_i$
  & 200--500  & \euro/day \\
$vhc_i$
  & 0.5--2.0  & \euro/m\textsuperscript{3} \\
\bottomrule
\end{tabular}
\end{table}

% ══════════════════════════════════════════════════════════════════
%  EXPERIMENTAL SETUP
% ══════════════════════════════════════════════════════════════════

\subsection{Experimental Setup}
\label{sec:exp-setup}

All experiments are run on a workstation equipped with an
Intel\textsuperscript{\textregistered} Core\texttrademark\
i7-XXXXX processor (X cores, X.X~GHz) and XX~GB RAM, running
Ubuntu~XX.XX. The linearised MILP is solved with Gurobi~XX.X
using default settings, with a time limit of 3\,600~s for
\textsc{Small} and \textsc{Medium} instances and 7\,200~s for
\textsc{Large} instances. The optimality gap tolerance is set to
0.5\,\%.

% ══════════════════════════════════════════════════════════════════
%  EXPERIMENT 1: SCALABILITY
% ══════════════════════════════════════════════════════════════════

\subsection{Experiment~1: Computational Scalability}
\label{sec:exp-scalability}

The first experiment evaluates how the solution time and the
optimality gap scale with the instance dimensions. For each
instance class, Table~\ref{tab:scalability} reports the number of
binary and continuous variables, the number of constraints, the
average solution time, the median optimality gap at termination,
and the fraction of instances solved to the gap tolerance within
the time limit.

\begin{table}[ht]
\centering
\caption{Computational scalability across instance classes.
  Times and gaps are averaged over $N_{\text{inst}}$ instances
  per class.}
\label{tab:scalability}
\small
\begin{tabular}{lrrrrrr}
\toprule
\textbf{Class}
  & \textbf{Bin.\ vars}
  & \textbf{Cont.\ vars}
  & \textbf{Constrs}
  & \textbf{Avg.\ time (s)}
  & \textbf{Med.\ gap (\%)}
  & \textbf{Solved (\%)} \\
\midrule
\textsc{Small}  & --- & --- & --- & --- & --- & --- \\
\textsc{Medium} & --- & --- & --- & --- & --- & --- \\
\textsc{Large}  & --- & --- & --- & --- & --- & --- \\
\bottomrule
\end{tabular}
\end{table}

To identify the dominant source of complexity, we additionally
report solve times as a function of each individual scaling
dimension---$|V|$, $|E|$, and $|R|$---while holding the other two
approximately constant.

% ══════════════════════════════════════════════════════════════════
%  EXPERIMENT 2: BASELINE COMPARISON
% ══════════════════════════════════════════════════════════════════

\subsection{Experiment~2: Cost Savings Relative to the One-to-One Baseline}
\label{sec:exp-baseline}

The second experiment compares the optimised SND-RTI solution with
the one-to-one baseline described in Section~\ref{sec:problem}. In
the baseline, every full-RTI route $(i,j)$ is paired with a mirror
return route $(j,i)$ that ships back the exact number of empty RTIs
consumed at~$j$, using the same transport mode. No consolidation or
rerouting of empty RTIs is allowed, and no hubs are activated.

For each instance, we compute the baseline cost by fixing the empty
return routes to their one-to-one configuration and solving the
remaining model (queue times and RTI pool) to optimality. The
relative savings of the SND-RTI solution are decomposed into three
components:

\begin{enumerate}
  \item \textbf{Transport cost savings}
        ($\Delta TC$): reduction in the total daily transport cost
        $\overline{TC}$, driven by route consolidation and mode
        switching on the return legs.
  \item \textbf{RTI pool savings}
        ($\Delta P$): reduction in the number of purchased RTIs,
        achieved by shorter return cycles and better pool
        utilisation.
  \item \textbf{Hub cost}
        ($\Delta CH$): the additional hub-handling cost incurred
        by the optimised solution, partially offsetting the above
        savings.
\end{enumerate}

Table~\ref{tab:baseline} reports the average and range of these
savings across all instances within each class.

\begin{table}[ht]
\centering
\caption{Cost savings of SND-RTI over the one-to-one baseline
  (average $\pm$ std.\ dev.\ across instances).}
\label{tab:baseline}
\small
\begin{tabular}{lrrrr}
\toprule
\textbf{Class}
  & $\Delta TC$ (\%)
  & $\Delta P$ (\%)
  & $\Delta CH$ (\euro/day)
  & \textbf{Total savings (\%)} \\
\midrule
\textsc{Small}  & --- & --- & --- & --- \\
\textsc{Medium} & --- & --- & --- & --- \\
\textsc{Large}  & --- & --- & --- & --- \\
\bottomrule
\end{tabular}
\end{table}

We further analyse how the network density---measured as the ratio
$|E|/|V|$---affects the magnitude of savings, since denser networks
offer more consolidation opportunities for empty-RTI returns.

% ══════════════════════════════════════════════════════════════════
%  DISCUSSION PLACEHOLDER
% ══════════════════════════════════════════════════════════════════

\subsection{Discussion}
\label{sec:exp-discussion}

% TODO: populate after running experiments.
% Expected structure:
% - Scalability: which dimension drives complexity (likely |E| x |R|)
% - Baseline: savings increase with network density and RTI type diversity
